\documentclass[11pt,a4paper]{article}
\usepackage[utf8]{inputenc}
\usepackage{amsmath,amssymb,amsfonts}
\usepackage{geometry}
\geometry{margin=1in}
\usepackage{booktabs}
\usepackage{hyperref}

\title{Exclusion and Visibility Geometry in VibeVolts}
\author{Simulation Reference Documentation}
\date{\today}

\begin{document}

\maketitle

\section{Introduction}
This document details the geometric, vector, and angular calculations used in VibeVolts to determine whether a target point in space is visible to a satellite or ground-based observatory detector. 

Visibility is subject to two conditions:
\begin{enumerate}
    \item \textbf{Field-of-View (FOV) Check}: The target must lie within the sensor's conical field-of-view.
    \item \textbf{Exclusion Checks}: The detector's view must not be obstructed by, or lie too close to, key celestial bodies (Sun, Moon, Earth).
\end{enumerate}

---

\section{Step 1: Coordinate Frame and Vectors}
Let $\mathbf{x}_{\text{sat}}$ be the position of the observer (satellite or observatory) and $\mathbf{x}_{\text{targ}}$ be the position of the target, both represented in the Geocentric Celestial Reference System (GCRS).

The pointing vector from the observer to the target is:
\begin{equation}
    \mathbf{v}_{\text{targ}} = \mathbf{x}_{\text{targ}} - \mathbf{x}_{\text{sat}} \quad [\text{meters}]
\end{equation}

The unit pointing vector is:
\begin{equation}
    \hat{\mathbf{v}}_{\text{targ}} = \frac{\mathbf{v}_{\text{targ}}}{\|\mathbf{v}_{\text{targ}}\|}
\end{equation}
If the distance $\|\mathbf{v}_{\text{targ}}\| < 10^{-9}$ meters, the target is considered co-located with the satellite and is automatically marked as excluded.

---

\section{Step 2: Field-of-View (FOV) Cone Check}
Let $\hat{\mathbf{p}}$ be the unit vector representing the pointing direction of the detector, and let $\text{FOV}$ be the full-width angular diameter of the field-of-view in radians.

The angle $\theta_{\text{targ}}$ between the detector pointing vector and the target direction is:
\begin{equation}
    \theta_{\text{targ}} = \arccos(\hat{\mathbf{p}} \cdot \hat{\mathbf{v}}_{\text{targ}}) \quad [\text{radians}]
\end{equation}

The target is within the detector's field of view if and only if:
\begin{equation}
    \theta_{\text{targ}} < \frac{\text{FOV}}{2}
\end{equation}

---

\section{Step 3: Apparent Angular Radii of Celestial Bodies}
As viewed from the observer's position $\mathbf{x}_{\text{sat}}$, celestial bodies appear as disks with angular sizes that depend on their physical radii and their distance from the observer.

Let $\mathbf{x}_{\text{body}}$ be the GCRS position of a celestial body (Sun, Moon, or Earth) and $R_{\text{body}}$ be its physical radius. The vector from the observer to the body is:
\begin{equation}
    \mathbf{v}_{\text{body}} = \mathbf{x}_{\text{body}} - \mathbf{x}_{\text{sat}}
\end{equation}
And the unit vector is $\hat{\mathbf{v}}_{\text{body}} = \mathbf{v}_{\text{body}} / \|\mathbf{v}_{\text{body}}\|$.

The apparent angular radius $\theta_{\text{body}}$ of the body as seen by the observer is:
\begin{equation}
    \theta_{\text{body}} = \arctan\left(\frac{R_{\text{body}}}{\|\mathbf{v}_{\text{body}}\|}\right) \quad [\text{radians}]
\end{equation}

In VibeVolts, the physical parameters and apparent radii calculations are:
\begin{itemize}
    \item \textbf{Sun:} Modeled as a point source for simplicity:
    \begin{equation}
        \theta_{\text{Sun}} = 0.0 \quad [\text{radians}]
    \end{equation}
    \item \textbf{Moon:} Using physical Moon radius $R_{\text{Moon}} \approx 1,737,400\text{ m}$:
    \begin{equation}
        \theta_{\text{Moon}} = \arctan\left(\frac{R_{\text{Moon}}}{\|\mathbf{v}_{\text{Moon}}\|}\right)
    \end{equation}
    \item \textbf{Earth:} Center is at GCRS origin $[0,0,0]$, so $\mathbf{v}_{\text{Earth}} = -\mathbf{x}_{\text{sat}}$. Using equatorial radius $R_{\text{Earth}} \approx 6,378,137\text{ m}$:
    \begin{equation}
        \theta_{\text{Earth}} = \arctan\left(\frac{R_{\text{Earth}}}{\|\mathbf{x}_{\text{sat}}\|}\right)
    \end{equation}
\end{itemize}

---

\section{Step 4: Exclusion Angle Check}
To protect sensitive optics from bright sources, a target observation is excluded if the line of sight passes too close to the limb (edge) of a celestial body.

Let $\psi_{\text{body}}$ be the angular separation between the target direction and the celestial body center:
\begin{equation}
    \psi_{\text{body}} = \arccos(\hat{\mathbf{v}}_{\text{targ}} \cdot \hat{\mathbf{v}}_{\text{body}}) \quad [\text{radians}]
\end{equation}

Let $\theta_{\text{limit}}$ be the exclusion safety angle specified for the detector (e.g. \texttt{solarEx}, \texttt{lunarEx}, \texttt{earthEx} in radians). The angle between the target line-of-sight and the limb of the body is $\psi_{\text{body}} - \theta_{\text{body}}$.

The view is **excluded** (obstructed or blinded) if:
\begin{equation}
    \psi_{\text{body}} - \theta_{\text{body}} < \theta_{\text{limit}}
\end{equation}

If this inequality holds for the Sun, Moon, or Earth, the target is excluded from observation.

---

\section{Summary of Variable Definitions}
\begin{table}[h]
\centering
\begin{tabular}{lll}
\toprule
Variable & Code Representation & Description \\
\midrule
$\mathbf{x}_{\text{sat}}$ & \texttt{sat\_pos} / \texttt{satposition} & Observer GCRS position vector (m) \\
$\mathbf{x}_{\text{targ}}$ & \texttt{targets} & Target GCRS position vector (m) \\
$\hat{\mathbf{p}}$ & \texttt{pointing} & Sensor pointing direction unit vector \\
$\text{FOV}$ & \texttt{fov} & Sensor full field-of-view (rad) \\
$\mathbf{v}_{\text{body}}$ & \texttt{vecs\_to\_bodies} & Vector from observer to celestial body center (m) \\
$\theta_{\text{body}}$ & \texttt{apparent\_radii} & Apparent angular radius of celestial body (rad) \\
$\theta_{\text{limit}}$ & \texttt{solarEx}, \texttt{lunarEx}, \texttt{earthEx} & Detector exclusion safety angles (rad) \\
$\psi_{\text{body}}$ & \texttt{angles} & Separation angle between target and body center (rad) \\
\bottomrule
\end{tabular}
\caption{Exclusion parameter and code variable mappings.}
\end{table}

\end{document}
